% !TeX root = ../thesis.tex
\chapter{Fazit und Ausblick}
\label{chap:fazit}

\section{Zusammenfassung der zentralen Ergebnisse}
\label{sec:fazit-zusammenfassung}

Ziel dieser Bachelorarbeit war es, \ac{ki}-basierte Ansätze zur automatisierten Informationsextraktion aus ausgefüllten Teilnehmerlisten im Kontext des Schulungswesens der BMW Group zu untersuchen, um künftig die Übertragung neu eingehender Teilnehmerlisten in das Schulungssystem zu automatisieren, und die Verfahren anhand des historischen Referenzbestands hinsichtlich ihrer Extraktionsqualität, ihrer Laufzeit- und Ressourceneffizienz sowie ihrer praktischen Einsatzfähigkeit systematisch zu bewerten. Hierzu wurde ein vollständig lokal unter einem Grafikspeicherbudget von \mbox{16\,GiB} \ac{vram} ausführbares, modulares Evaluationsframework in Python~3.12 entwickelt, das fünf ausgewählte Extraktionsverfahren (\texttt{ocr\_rules}, \texttt{ocr\_llm}, \texttt{paddle\_llm}, \texttt{vlm}, \texttt{ocr\_vlm} in \texttt{start\_testframework\_5.cmd} / \texttt{configs/benchmark-five.yaml}) unter kontrollierten, reproduzierbaren Bedingungen vergleicht.

Die zentralen wissenschaftlichen und technischen Ergebnisse der Arbeit lassen sich in fünf Punkten zusammenfassen:
\begin{enumerate}
  \item \textbf{Überlegenheit strukturierter \ac{ocr}+\ac{llm}- und hybrider \ac{ocr}+\ac{vlm}-Pipelines:} Auf den untersuchten realen Entwicklungsstichproben erzielen \texttt{paddle\_llm} (PP-StructureV3/PP-OCRv5 + Text-\ac{llm} \texttt{Qwen3-4B-Instruct-2507}; $51{,}09\,\%$ exakter Feld-Recall, $60{,}65\,\%$ exakter Feld-$F_1$ und $68{,}48\,\%$ Feldabdeckung im Dreidokumentlauf) sowie das hybride Verfahren \texttt{ocr\_vlm} (vorgeschaltete PaddleOCR kombiniert mit \texttt{Qwen3-VL-4B-Instruct}; $58{,}33\,\%$ exakter Feld-Recall auf dem gemischten Dokument \texttt{140829.pdf} und $52{,}17\,\%$ im Dreidokumentlauf nach abgesicherter Folgeseitenreparatur) die höchste Extraktionsgüte. Sie übertreffen sowohl die klassische Tesseract-\ac{llm}-Kombination (\texttt{ocr\_llm}: $40{,}22\,\%$) als auch das rein bildbasierte \ac{vlm} (\texttt{vlm}: $36{,}41\,\%$) und die vorlagengebundene Regelbaseline (\texttt{ocr\_rules}: $0{,}00\,\%$) deutlich.
  \item \textbf{Entscheidende Rolle vor- und nachgelagerter Pipeline-Mechanismen:} Die kausale Fehleranalyse über bis zu $20$~Seiten ($620$~Referenzfelder) hat gezeigt, dass massive Feldverluste in realen Dokumentarchiven häufig durch strukturelle Brüche entlang der Pipeline entstehen. Durch die Entwicklung (a)~einer zweistufigen automatischen Seitenorientierung (\ac{osd} + Vierlagen-\ac{ocr}-Fallback mit $20/20$ korrekt ausgerichteten Entwicklungsseiten), (b)~einer zerstörungsfreien Hintergrundnormalisierung für grau hinterlegte Kopfzeilen, (c)~eines verlustfreien \ac{ocr}-Beobachtungsregisters, (d)~einer geometrisch geprüften Spaltenvererbung für kopflose Folgeseiten ($\Delta x_{\max} \le 0{,}035$) mit nativer Anker-Absicherung (\texttt{echo\_corroboration}), (e)~einer zeilenindividuellen Vor- und Nachname-Trennung sowie (f)~einer verlustfreien Input-Kompaktierung gegen Tokenbudget-Abbrüche konnten diese strukturellen Fehlerquellen gezielt behoben werden.
  \item \textbf{Methodische Präzisierung der Evaluationsmetriken:} Durch die Trennung von primärer Originalrollenbewertung (\texttt{original\_roles}) und Namensswap-Sensitivitäten, die Einführung eines strikten Teilnehmer-Recalls, die fünfstufige Feldstatus-Partitionierung (\emph{Korrekt}, \emph{Abweichend}, \emph{Nicht erkannt}, \emph{Zuordnung unklar}, \emph{Nicht bewertbar}) und den Ersatz des mehrdeutigen Soft-$F_1$ durch die bedingte mittlere Levenshtein-Ähnlichkeit erkannter Felder ($S_{\mathrm{detected}} = \Sigma_s / |D|$) neben der Feldabdeckung ermöglicht das Framework eine unverzerrte, bis auf jedes Einzelfeld prüfbare Qualitätsdiagnose.
  \item \textbf{Ressourceneffizienz unter lokaler Hardwaregrenze:} Mithilfe isolierter \texttt{uv}-Laufzeitumgebungen, sequenzieller Prozessentladung und \ac{bf16}-Inferenz arbeiten sämtliche fünf Verfahren zuverlässig unterhalb von \mbox{16\,GiB} \ac{vram} (Spitzenallokation $\le 10\,\text{GiB}$). Zudem verhindert die Kombination aus PaddleOCR-Hinweisen und \ac{vlm} in \texttt{ocr\_vlm} auf dem Dreidokumentbestand die beim reinen \ac{vlm} beobachteten langwierigen Ausgabetoken-Abbrüche.
  \item \textbf{Datenschutzkonforme Bereitstellung:} Die strikte Offline-Ausführung des lokalen Frameworks und die Bereitstellung eines quellcodegleichen, datenfreien Transferpakets (\texttt{software/api\_firmen\_pc/}) gewährleisten die Einhaltung betrieblicher Datenschutzvorgaben und wurden durch insgesamt $922$ lokale Tests und $48$ \ac{api}-Tests verifiziert.
\end{enumerate}

\section{Ausblick}
\label{sec:fazit-ausblick}

Aufbauend auf dem geschaffenen Evaluationsframework ergeben sich fünf konkrete Ansatzpunkte für weiterführende wissenschaftliche und praktische Arbeiten:
\begin{itemize}
  \item \textbf{Durchführung der konfirmatorischen Pilot- und Finalstudie:} Als unmittelbarer nächster Schritt ist aus dem historischen Referenzbestand von rund $6\,000$ bereits erfassten Dokumenten über den implementierten Auswahlworkflow (\texttt{participant-bench selection inventory/freeze/run/compare}) die geschichtete Pilotstichprobe von $30$ bis $50$ unabhängigen Listen sowie anschließend die gesperrte finale Testevaluation auf mindestens zehn unabhängigen Schulungsgruppen zu ziehen und auszuwerten, um belastbare 95-\%-Cluster-Bootstrap-Konfidenzintervalle zu bestimmen.
  \item \textbf{Unabhängige Seiten- und Leerzellenannotation:} Um die Diskrepanz zwischen externen Excel-Kurslisten und den tatsächlich auf dem Papierformular sichtbaren Spalten aufzulösen, sollte für die Teststichprobe über \texttt{ocr-reference} eine unabhängige seitenbezogene Referenzannotation (einschließlich expliziter \texttt{empty}-Zellen, \texttt{signature\_present}-Labels und Seitentranskripten für Handschrift-\ac{cer}/\ac{wer}) erstellt und die mehrstufige Konfidenzkalibrierung (\texttt{ocr-calibrate} / \texttt{ocr-final}) abgeschlossen werden.
  \item \textbf{Formale Grammatikbeschränkung gegen \ac{vlm}-Tokenabbrüche:} Um die im direkten und hybriden \ac{vlm}-Pfad auf dichten Folgeseiten beobachteten Abbrüche am $8\,192$-Token-Ausgabelimit zu verhindern, sollte eine formale Grammatik- oder Zustandsautomaten-Beschränkung während der autoregressiven Dekodierung~\cite{Geng2023Grammar} sowie eine kontrollierte Zeilensegmentierung bei sehr langen Tabellen untersucht werden.
  \item \textbf{Evaluation weiterer lokaler Modellgenerationen und spezialisierter \ac{htr}-Zeilenmodelle:} Die modulare Worker-Architektur erlaubt es, künftig neuere Erkennungs- und Sprachmodelle (wie \texttt{PP-OCRv6}, \texttt{PaddleOCR-VL-1.6} oder quantisierte \texttt{Qwen3-8B}-Varianten) sowie auf den durch PP-StructureV3 lokalisierten Zellausschnitten spezialisierte Handschrift-Transkriptionsmodelle wie feinabgestimmtes \ac{trocr}~\cite{Li2023TrOCR} ohne Änderung der Evaluationslogik zu integrieren.
  \item \textbf{Produktive Human-in-the-Loop-Integration für künftige Schulungslisten:} Schließlich kann das erzeugte Audit- und Dashboard-Format (\texttt{report.html}, \texttt{participants.xlsx}) über eine Freigabeschnittstelle direkt mit dem betrieblichen Schulungsmanagement-System gekoppelt werden, sodass bei künftig eingehenden Teilnehmerlisten die vom Prüfpersonal bestätigten oder korrigierten Teilnehmerdatensätze automatisiert und ohne erneute manuelle Eingabe übernommen werden.
\end{itemize}
